\chapter{Objetivos}

\section{Geral}

Desenvolver o DCOU, um \textit{software} \textit{web} completo e modular para cálculo e dimensionamento de equipamentos de operações unitárias, integrando rotinas de balanço de massa, transferência de calor e reações químicas e, simultaneamente, oferecendo recursos didáticos que apoiem a compreensão dos conceitos envolvidos. O sistema deve adotar arquitetura cliente-servidor com interface \textit{web} acessível por meio de API aberta, mantendo o código disponível publicamente.

\section{Específico}

\begin{itemize}
    \item Implementar rotinas de cálculo baseadas em equações clássicas para operações unitárias como reatores, tubulações e demais correlatos, com perda de carga, \textit{NPSH} etc;
    \item Desenvolver uma interface \textit{web} interativa e responsiva para inserção e interpretação de dados em diversos dispositivos, favorecendo a experiência do usuário;
    \item Integrar um banco de dados local (inicialmente \textit{hardcoded}) com propriedades físico-químicas de substâncias comuns (coeficientes de difusão, calor específico, entalpia de vaporização etc.);
    \item Disponibilizar recursos didáticos integrados --- glossário técnico, explicações teóricas embutidas e visualizações interativas dos fenômenos físicos --- que conectem cada cálculo ao seu fundamento conceitual;
    \item Desenvolver exercícios guiados que encadeiem múltiplos módulos, reproduzindo fluxos de trabalho reais de engenharia;
    \item Estruturar trilhas de aprendizagem que organizem o uso dos módulos em sequência didática;
    \item Validar os resultados comparando-os com dados da literatura, assegurando precisão e confiabilidade dos cálculos.
\end{itemize}
